
As our running example, we will consider the engine type $\mathsf{calc}$ that
can process messages of type in \eqref{eq:message-interface-example}. The engine
type $\mathsf{calc}$ would be defined later, once all required constructs are
presented.

% The following is our running example of an engine.
\begin{example}[Calculator Engine]\label{ex:calculator-engine}%
  Let us assume the system defines the engine type $\enginetype{\mathsf{calc}}$
  and that an engine of this type works as normal calculator with some internal
  memory. The engine has the following message interface:
  \begin{equation}\label{eq:message-interface-example}
  \typedef%
  {\msgtype{\mathsf{calc}}}%
  {    \msg{Add}(\mathsf{Pair}) 
  \mid \msg{Subtract}(\mathsf{Pair}) 
  \mid \msg{Multiply}(\mathsf{Pair})
  \mid \msg{Result}(\mathsf{Nat})
  },
\end{equation}

where $\msg{Add}$, $\msg{Subtract}$, and $\msg{Multiply}$ are message tags, representing
the operations that the calculator engine can perform to modify its memory. The
message type $\mathsf{Pair}$ these message tags wrap encodes two natural numbers,
representing the two operands of the operation. The message tag $\msg{Result}$ wraps
a natural number, representing the result stored in the engine's memory.
\end{example}